\scanpage{12}
\mspage{10}
Pour en déduire une application sur les $\varinjlim$
\[
\underbrace{\operatorname{Hom}_\Sigma(x,y)}_{\substack{\Vert\\\varinjlim\operatorname{Hom}_{\Sigma,n}(x,y)}}
\times\underbrace{\operatorname{Hom}_\Sigma(y,z)}_{\substack{\Vert\\\underline{\lim}\operatorname{Hom}_{\Sigma,m}(y,z)}}
\longrightarrow\underbrace{\operatorname{Hom}_{\Sigma\deleted{\unsure{+}}}(x,z)}_{\substack{\Vert\\\underline{\lim}\operatorname{Hom}_{\Sigma,\unsure{n+m}}(x,z)}},
\]
il faut prouver \deleted{la}, pour $n\leq n',m\leq m'$, la commutativité de
\diagram{\begin{tikzcd}[ampersand replacement=\&,column sep=.8em,scale=.82,transform shape]
\operatorname{Hom}_{\Sigma,\unsure{n|m}}(x,y)\times\operatorname{Hom}_{\Sigma,\unsure{m|n}}(y,z)\arrow[r]\arrow[d] \& \operatorname{Hom}_{\Sigma,m+n}(x,y)\arrow[d]\\
\operatorname{Hom}_{\Sigma,n'}(x,y)\times\operatorname{Hom}_{\Sigma,m'}(y,z)\arrow[r] \& \operatorname{Hom}_{\Sigma,m'+n'}(x,y).
\end{tikzcd}}

En d'autres termes, que pour $c:x\longrightarrow \unsure{y'|y}$ et $c':y\longrightarrow z$ deux \inserted{$\Sigma$-}chemins composables de long. $n$ et $m$ respectivement, le composé \deleted{$\unsure{c'c}$} $c'c:x\longrightarrow z$ ne change pas, ``à homotopie près'' dans $\underline{\mathrm{Ch}}_\Sigma(x,z)$, quand on remplace $c,c'$ par les chemins $\bar c,\bar c'$ de longueurs $\unsure{n',m'|m',n'}$ déduits de $c$ et $c'$ par ``dégénérescences aux bouts''.

\deleted{Ainsi, \unsure{on} \unsure{voit que le chemin}} \inserted{Plus précisément,}
\[
\overline{c'}\,\bar c\ \deleted{\unsure{\sim}}\ \overline{c'c},
\]
le chemin composé est homotope à celui déduit de $c'c$ \unsure{par} ``dégénérescences aux bouts''.

Je viens de regarder pour $n=m=1$, $n'\deleted{\unsure{+}}=m'=\unsure{2|3|?}$, et de constater qu'il n'y a pas de \unsure{flèche|flèches} \unsure{naturelle|naturelles|reliant} entre $\overline{c'}\,\bar c$ et $\overline{c'c}$ -- l'homotopie cherchée doit

% archive sheet number: 11
